\section{Research Gaps}

Despite significant progress in each of the areas reviewed above,
several gaps remain. First, pruning, quantisation, and distillation are
typically applied as independent post-processing steps, ignoring
potentially synergistic interactions between them. Second, most existing
methods optimise for a single efficiency metric (e.g., FLOPs) rather
than jointly considering the multiple resource dimensions---compute,
memory, and energy---that matter in practice. Third, hardware-aware
approaches tend to be tightly coupled to specific device
characteristics, limiting their generality.

The RANO framework proposed in this thesis addresses these gaps by
providing a unified, multi-objective optimisation procedure that
jointly optimises architecture, sparsity, and precision while remaining
adaptable to diverse hardware targets through a pluggable cost model.

\begin{table}[htbp]
\centering
\caption{Comparison of model compression techniques. ``Unstructured''
refers to fine-grained weight removal; ``Structured'' refers to channel
or filter removal. Accuracy figures are representative for ResNet-50 on
ImageNet.}
\label{tab:compression}
\begin{tabular}{@{}lccc@{}}
\toprule
\textbf{Method} & \textbf{Compression} & \textbf{Top-1 Acc.} & \textbf{Speedup} \\
\midrule
Unstructured pruning    & 10$\times$ & 74.2\% & 1.2$\times$ \\
Structured pruning      & 3$\times$  & 75.8\% & 2.8$\times$ \\
INT8 quantisation       & 4$\times$  & 76.0\% & 3.5$\times$ \\
Knowledge distillation  & 5$\times$  & 75.3\% & 4.0$\times$ \\
Joint (RANO, ours)      & 8$\times$  & 76.1\% & 5.2$\times$ \\
\bottomrule
\end{tabular}
\end{table}

\newpage

% ═════════════════════════════════════════════════════════════════════════════
% CHAPTER 3
% ═════════════════════════════════════════════════════════════════════════════
